\begin{otherlanguage}{english}
\chapter*{Abstract}
\addcontentsline{toc}{chapter}{Abstract}

In Benin, a growing number of businesses, schools, administrations and SMEs run a
website without the skills or resources needed to verify whether it meets basic
security standards. Existing tools (Mozilla Observatory, SSL Labs) are often
English-only, scattered, and not designed for non-technical users. As part of our
academic internship, we designed and developed \textbf{CyberGuard B\'enin}, a web
platform that automates this security diagnostic: the user enters a URL, and the
system runs a series of non-intrusive checks (HTTPS, SSL certificate, security
headers, cookies, server information disclosure), then translates these results
into a score out of 100, a risk level, and clear, actionable recommendations. The
project began with a clarification of the need and objectives, followed by a UML
modeling phase (actors, use cases, class, object and sequence diagrams) and
database design. The technical implementation relies on a Laravel (PHP) API on
the backend and a React application on the frontend. This project offers a
concrete solution to democratize access to a first level of security auditing for
website owners in Benin.

\bigskip
\noindent\textbf{Keywords:} web security, automated audit, HTTPS, Laravel, React,
UML.
\end{otherlanguage}
